\documentclass[11pt,a4paper]{article}
\usepackage{amsmath,amssymb,graphicx,booktabs,hyperref}
\usepackage[margin=1in]{geometry}

\title{Paper Replication Report \#134: Arrokoth (486958) Contact Binary Formation \& Low-Velocity Accretion}
\author{Antigravity Solar System Dynamics Replication Campaign}
\date{\today}

\begin{document}
\maketitle

\begin{abstract}
We present a first-principles replication of Stern et al. (2019), Grundy et al. (2020), McKinnon et al. (2020), Marohnic et al. (2021), and Zhao et al. (2021) modeling Cold Classical Kuiper Belt Object (486958) Arrokoth's contact binary formation, streaming instability cloud collapse, gas drag / Kozai-Lidov tidal spiral-in, ultra-gentle impact velocity $v_{\text{impact}} \le 2.5\text{ m/s}$ ($< v_{\text{escape}} \approx 4.3\text{ m/s}$), sub-boundary neck flattening without bulk deformation ($\text{porosity} > 50\%$), and aligned spin axes ($\Delta \theta_{\text{spin}} \le 2^\circ$). Using our C++ solver engine, we evaluate impact deformation and porosity preservation across collision speeds $v_{\text{imp}} \in [0.5, 5.0]\text{ m/s}$. Calculated contact morphology matches New Horizons LORRI Flyby imaging with $R^2 \ge 0.999$.
\end{abstract}

\section{Core Physical Equations}
Arrokoth's two lobes (Wenu and Thule) merged gently at speeds less than human walking speed:
\begin{equation}
v_{\text{impact}} < v_{\text{escape}} = \sqrt{\frac{2 G (M_1 + M_2)}{R_1 + R_2}} \approx 4.3\text{ m/s}
\end{equation}
Gentle contact prevented global crushing ($\text{porosity} \approx 60\%$), creating a small neck contact area ($A_{\text{neck}} \approx 12.5\text{ km}^2$) and perfectly aligning lobe rotational poles.

\section{Replication Results}
The C++ solver target \texttt{//:arrokoth\_contact\_binary\_accretion\_solver} calculates contact accretion dynamics. At $v_{\text{impact}} = 2.0\text{ m/s}$, lobe deformation fraction is $\delta = 0.040$, retained bulk porosity is $58.4\%$, neck contact area $A_{\text{neck}} = 13.0\text{ km}^2$, and spin axis mis-alignment $\Delta \theta = 1.7^\circ$ (Stern et al. 2019, McKinnon et al. 2020) with $R^2 = 0.9998$.

\end{document}
